\chapter{Literature Review}
\label{chap:literature}

\section{Overview}

Topic segmentation has progressed through three broad generations.
\emph{Classical} lexical-cohesion methods from the 1990s and 2000s detected
boundary positions where local vocabulary shifts, using sliding-window
similarity scores on bag-of-words representations~\citep{hearst1994multi,choi2000latent}.
\emph{Neural} methods from 2018 onwards replaced sparse lexical features with
dense contextual sentence embeddings, improving performance on diverse genres
at the cost of requiring pretrained language models~\citep{koshorek2018text,reimers2019sbert}.
\emph{Multimodal} methods subsequently incorporated visual, audio, and structural
signals from videos and meeting recordings, but often target different
benchmarks, larger supervised settings, or task-specific chapter-generation
metrics~\citep{yu2024_multimodal,cao2022_chaptergen,yang2023_vidchapters}. This thesis therefore treats prior work as
context rather than as a directly comparable leaderboard: LECSEG is positioned
as a low-resource, reproducible lecture benchmark and pipeline.

\section{Classical text-based segmentation}
\label{sec:classical}

\subsection{TextTiling}

TextTiling~\citep{hearst1994multi} partitions expository text by computing
pairwise lexical similarity between adjacent blocks of $k$ sentences,
using a sliding window of width $w$.
Boundaries are placed at local minima (``valleys'') of the similarity profile
after smoothing.
The method is parameter-sensitive (window size $w$, block size $k$) and
degrades when technical vocabulary is dense and consistent throughout a lecture.
Our implementation adds adaptive defaults: $w = \min(w, \max(3, N/10))$ and
$k = \min(k, \max(2, n_{\text{blocks}}/3))$, which improve performance on
the short segments typical of lecture videos.

\subsection{C99 and rank-transformed cosine similarity}

C99~\citep{choi2000latent} computes a cosine-similarity matrix over
bag-of-words sentence vectors, applies a rank transformation to compress
the dynamic range, and uses a diagonally dominated clustering criterion to
locate boundaries.
It is more robust than TextTiling on technical lectures because the rank
transformation reduces sensitivity to term frequency outliers.
We use C99 as one of our classical baselines (B2).

\subsection{BayesSeg and topic models}

Eisenstein and Barzilay~\citep{eisenstein2008bayesian} formulated segmentation
as inference in a hierarchical Bayesian topic model, naturally producing
uncertainty estimates over boundaries.
While principled, these models are slow and require careful hyperparameter
tuning; they are omitted from our main comparison but included in related work.

\section{Neural sentence-embedding methods}
\label{sec:neural}

\subsection{Cosine-drop with contextual embeddings}

The advent of sentence transformers~\citep{reimers2019sbert} enabled a simple
but effective baseline: compute the cosine similarity between adjacent sentence
embeddings, and threshold the depth-score valleys.
When applied to SBERT or MPNet embeddings of lecture transcripts, this approach
significantly outperforms TextTiling and C99 with no training, motivating its
inclusion as baseline B3.

\subsection{Supervised neural segmentation}

Koshorek et al.~\citep{koshorek2018text} trained the first supervised neural
segmenter (TextSeg) using BERT sentence embeddings with a classification head.
Li et al.~\citep{li2018segbot} proposed SegBot, which uses a biLSTM to score
potential boundaries.
Both require domain-matched labelled data; on out-of-domain lecture content,
zero-shot cosine baselines often match or exceed them.
We use the cosine-depth variant (BertSeg, B5) as the strongest neural baseline.

\section{Multimodal lecture and meeting segmentation}
\label{sec:multimodal_lit}

\subsection{Meeting and conversation segmentation}

Meeting and conversation segmentation methods often combine lexical features,
speaker turns, and discourse markers. These approaches target short-form
conversational content and do not directly solve hour-long academic monologues.

\subsection{Lecture video segmentation}

Recent multimodal lecture-video segmentation work uses text and visual features
with supervised or transformer-based fusion, achieving strong results when
labelled training data are available~\citep{yu2024_multimodal}.
Such systems are important reference points, but their datasets and training
assumptions differ from the low-resource LECSEG setting.
Slide-change segmentation methods use slide content changes alone,
which fails for whiteboard-style lectures with no slide transitions.

\lecseg{} differs from many prior multimodal lecture systems by emphasising
local reproducibility, shared metrics, and explicit negative results when OCR,
shot, or prosody cues fail to improve $P_k$/WindowDiff.

\subsection{Large-scale video chaptering benchmarks}

Recent video chaptering work has moved toward much larger supervised datasets.
Chapter-Gen localises and titles chapters in thousands of user-generated
videos~\citep{cao2022_chaptergen}. VidChapters-7M scales this direction to
hundreds of thousands of videos and millions of chapters, evaluating generation
and localisation with video-captioning and temporal-localisation metrics rather
than the $P_k$/WindowDiff pair used in text segmentation~\citep{yang2023_vidchapters}.
MiniSeg/YTSEG reports strong supervised smart-chaptering results on 19,299
YouTube videos~\citep{retkowski2024_miniseg}, and Chapter-Llama shows that
trained long-context LLM chaptering can be highly effective on VidChapters-scale
benchmarks~\citep{ventura2025_chapterllama}. TreeSeg is a closer small-data
transcript comparator because it reports $P_k$/WindowDiff on lecture and
meeting transcripts~\citep{gklezakos2024_treeseg}.

These systems are stronger than \lecseg{} in raw scale and, in several cases,
reported benchmark performance. \lecseg{} therefore should be positioned in a
different niche: it is a small, inspectable, lecture-specific benchmark and
reproducible pipeline that reports classical segmentation metrics,
hierarchical annotation, inter-annotator agreement, and failure analysis.
Table~\ref{tab:external_scale} summarises this scale difference, while
Table~\ref{tab:related_work_detailed} records the more detailed comparison
used to keep the thesis claim boundary explicit.
Table~\ref{tab:low_resource_positioning} then quantifies the narrow
low-resource claim: LECSEG uses a tiny fraction of the video count used by the
largest supervised chaptering systems, but this is a scale claim rather than an
external performance claim.

\begin{table}[t]
\centering
\caption{Video-count scale comparison for related lecture/video chaptering work.}
\label{tab:external_scale}
\resizebox{\linewidth}{!}{%
\begin{tabular}{p{0.20\linewidth}r p{0.17\linewidth}p{0.22\linewidth}p{0.23\linewidth}}
\toprule
Work & Videos & Scale & Metrics & Positioning \\
\midrule
LECSEG-30 & 30 & 32.52 h & Pk/WD/BS/F1@2 & Low-resource lecture benchmark \\
TreeSeg TinyRec & 21 & n/a & Pk/WD & Closest small transcript comparator \\
Videoaula & 34 & n/a & F1 / hierarchy & Lecture ToC corpus \\
LectureDE & 96 & n/a & F1 / hierarchy & German lecture corpus \\
AVLectures & 2,350+ & large STEM lectures & task-specific & Multimodal lecture resource \\
Chapter-Gen & 9,631 & user videos & AP/Recall/ROUGE & Supervised chapter generation \\
MiniSeg/YTSEG & 19,299 & 6,533 h & Pk/BS/F1 & Large supervised YouTube benchmark \\
VidChapters-7M & 817,000 & 7M chapters & SODA/localization & Large-scale video chaptering \\
\bottomrule
\end{tabular}
}
\end{table}


\begin{table}[t]
\centering
\scriptsize
\setlength{\tabcolsep}{3pt}
\renewcommand{\arraystretch}{1.08}
\caption{Detailed comparison with related lecture/video chaptering work. Metrics are not directly interchangeable across datasets.}
\label{tab:related_work_detailed}
\begin{tabularx}{\linewidth}{p{0.17\linewidth}p{0.10\linewidth}p{0.18\linewidth}X p{0.22\linewidth}}
\toprule
Work & Videos & Supervision & Metric/result & LECSEG verdict \\
\midrule
LECSEG balanced selector & 30 & LOO method selector & Pk=0.3588; WD=0.3739; F1@2=0.0893 & Reference low-resource result. \\
MiniSeg/YTSEG & 19,299 & Supervised & Pk=28.73; BS=35.74; F1=43.37 & External method stronger on scale/Pk. \\
Chapter-Gen & 9,631 & Supervised multimodal & AP=43.3; R@5s=76.1 & Stronger supervised localization. \\
VidChapters-7M & 817,000 & Large-scale finetuned & 7M chapters; SODA\_c=11.4 & Far stronger scale. \\
Chapter-Llama & 10k/8.1k & Trained LLM & F1=45.3 vs Vid2Seq 26.7 & Stronger trained chapterer. \\
TreeSeg/TinyRec & 21 & Unsupervised & TinyRec Pk=0.367 & Closest small comparator; no shared-run win. \\
AVLectures & 2,350+ & Self-supervised AV & 86 STEM courses & Stronger lecture-video scale. \\
\bottomrule
\end{tabularx}
\end{table}

\begin{table}[t]
\centering
\small
\caption{Low-resource scale positioning. Ratios compare video counts only; metrics are not directly interchangeable across datasets.}
\label{tab:low_resource_positioning}
\begin{tabularx}{\linewidth}{lrrX}
\toprule
Work & Videos & Scale & Interpretation \\
\midrule
TreeSeg / TinyRec & 21 & 0.7x & Closest small unsupervised transcript comparator; not a scale advantage for LECSEG. \\
Chapter-Gen & 9,631 & 321.0x & Supervised multimodal chapter generation; over 300x LECSEG video count. \\
MiniSeg / YTSEG & 19,299 & 643.3x & Supervised smart-chaptering benchmark; over 600x LECSEG video count. \\
AVLectures & 2,350 & 78.3x & Large lecture-video representation resource; roughly 78x LECSEG video count. \\
VidChapters-7M & 817,000 & 27,233.3x & Large-scale video chaptering benchmark; over 27,000x LECSEG video count. \\
Chapter-Llama & 10,000 & 333.3x & Trained on about 10k videos and evaluated on 8.1k test videos. \\
\bottomrule
\end{tabularx}
\end{table}

\section{Hierarchical and discourse-based approaches}
\label{sec:hierarchical_lit}

Hierarchical topic segmentation has been studied in the NLP community under the
name \emph{discourse parsing}~\citep{mann1988rst}.
Hierarchical variants of topic segmentation motivate keeping the
chapter/subtopic nesting constraint explicit in LECSEG, even though raw
boundary performance remains the primary evaluation challenge.

\section{Local open language models for NLP refinement}
\label{sec:llm_lit}

The release of Llama-3.1~\citep{meta2024llama3} and other open models has made
it feasible to run LLM inference locally at acceptable cost.
Comparative evaluations show that compact open-weight models can be strong on
many zero-shot NLP tasks~\citep{jiang2024mixtral}.
We use Llama-3.1-8B (GGUF Q4 via Ollama) for boundary refinement and title
generation, eliminating the reproducibility and privacy concerns of closed APIs.

\section{Gap analysis}
\label{sec:gap}

Table~\ref{tab:litmatrix} summarises the gaps in prior work that motivate
the contributions of this thesis. The key gap is not that no stronger video
chaptering systems exist; several large supervised systems are stronger at
scale. The gap is that low-resource lecture segmentation still lacks compact,
auditable benchmarks with chapter/subtopic structure, a unified metric suite,
and explicit reporting of what fails. \lecseg{} is designed for this
reproducible lecture-specific niche.

\begin{table}[t]
  \centering
  \caption{Literature matrix. \checkmark: addressed; (p): partial; ---: not addressed.}
  \label{tab:litmatrix}
  \small
  \begin{tabular}{lcccc}
    \toprule
    Method & Open & Multimodal & Hierarchical & Public benchmark \\
    \midrule
    TextTiling~\citep{hearst1994multi}       & \checkmark & --- & --- & --- \\
    C99~\citep{choi2000latent}               & \checkmark & --- & --- & --- \\
    BertSeg~\citep{koshorek2018text}         & \checkmark & --- & --- & (p) \\
    Multimodal VTS~\citep{yu2024_multimodal} & --- & \checkmark & --- & --- \\
    Meeting/topic hierarchy methods          & (p) & (p) & (p) & --- \\
    \midrule
    \lecseg{} (this work)                   & \checkmark & (p) & \checkmark & \checkmark \\
    \bottomrule
  \end{tabular}
\end{table}

\section{Synthesis for this thesis}
\label{sec:literature_synthesis}

The literature leads to four design choices in this thesis. First, classical
segmentation work justifies using $P_k$ and WindowDiff, but also shows why
lexical methods alone are fragile on technical lectures with repeated
vocabulary. Second, neural sentence embeddings motivate the BGE/E5 cross-model
methods because they provide strong zero-shot semantic signals without requiring
large labelled training data. Third, large video-chaptering systems demonstrate
that stronger results are possible at scale, but their data volume, supervision,
and metrics are not directly comparable with a 30-video lecture benchmark.
Fourth, multimodal and LLM work suggests useful auxiliary signals, but the
LECSEG results show that such signals must be calibrated to creator-level
chapter granularity; otherwise they over-segment at the subtopic or rhetorical
level. The novelty of \lecseg{} is therefore not a single new state-of-the-art
model, but a compact lecture-specific benchmark, unified evaluation protocol,
annotation reliability report, and diagnostic account of what low-resource
methods can and cannot do.
